% from ega4-16-fr.tex
\oldpage[IV-4]{32}
et proposons-nous de chercher s'il existe des \emph{$S$-morphismes}
$u:Y\to X$ tels que $u_0=u\circ j$ (autrement dit, s'il est possible de
compléter le diagramme précédent par la flèche pointillée $u$, de façon à le
laisser \emph{commutatif}).

Pour cela, considérons un ouvert affine $U=\operatorname{Spec}(C)$ de $Y$;
son image réciproque $j^{-1}(U)$ est l'ouvert affine
$U_0=\operatorname{Spec}(C/\mathfrak{L})$, où
$\mathfrak{L}=\Gamma(U,\mathscr{J})$, idéal \emph{de carré nul} dans $C$;
nous supposerons $U$ assez petit pour que $u_0(U_0)$ soit contenu dans un
ouvert affine $V=\operatorname{Spec}(B)$ de $X$, et
$g(U)=f(u_0(U_0))$ contenu dans un ouvert affine
$W=\operatorname{Spec}(A)$ de $S$, de sorte que $B$ et $C$ sont des
$A$-algèbres et que $u_0|U_0$ correspond à un $A$-homomorphisme $\psi$ de
$B$ dans $C/\mathfrak{L}$; soit $P(U_0)$ l'ensemble des restrictions $u|U$
des homomorphismes cherchés, qui correspondent canoniquement aux
$A$-homomorphismes d'algèbres $\varphi:B\to C$ tels que le \emph{composé}
\[
  B\xrightarrow{\varphi}C\longrightarrow C/\mathfrak{L}
\]
\emph{soit égal à $\psi$}. On sait donc (0, 20.1.1) que l'ensemble de ces
homomorphismes est vide ou de la forme
$\varphi_1+\operatorname{D\acute{e}r}_A(B,\mathfrak{L})$; lorsque $P(U_0)$ n'est pas
vide, le groupe additif $\operatorname{D\acute{e}r}_A(B,\mathfrak{L})$ opère par
addition dans $P(U_0)$, qui est donc un \emph{espace affine} pour le groupe
additif $\operatorname{D\acute{e}r}_A(B,\mathfrak{L})$ (ou encore un
\emph{espace homogène principal} (ou \emph{torseur}) \emph{sous}
$\operatorname{D\acute{e}r}_A(B,\mathfrak{L})$).

Remarquons maintenant que, puisque $\mathfrak{L}$ est muni d'une structure de
$B$-module au moyen de $\psi$, on a un \emph{isomorphisme}
$v\mapsto v\circ d_{B/A}$ de
$\operatorname{Hom}_B(\Omega_{B/A}^1,\mathfrak{L})$ sur
$\operatorname{D\acute{e}r}_A(B,\mathfrak{L})$ (0, 20.4.8). D'ailleurs, comme
$\mathfrak{L}$ est de carré nul, donc un $(C/\mathfrak{L})$-module, tout
$B$-homomorphisme $v:\Omega_{B/A}^1\to\mathfrak{L}$ peut être considéré
comme un $(C/\mathfrak{L})$-homomorphisme
$\Omega_{B/A}^1\otimes_B(C/\mathfrak{L})\to\mathfrak{L}$. Comme
$\mathscr{J}$ est de carré nul, il peut être considéré comme un
$\mathcal{O}_{Y_0}$-Module quasi-cohérent; introduisons le
$\mathcal{O}_{Y_0}$-Module
\[
  \mathscr{G}=\mathscr{H}om(u_0^*(\Omega_{X/S}^1),\mathscr{J});
  \tag{16.5.14.2}
  \label{IV.16.5.14.2}
\]
il résulte alors du fait que
$\Omega_{B/A}^1=\Gamma(V,\Omega_{X/S}^1)$ (16.3.7) que l'on peut écrire
$\operatorname{D\acute{e}r}_A(B,\mathfrak{L})=\Gamma(U_0,\mathscr{G})$.

Comme $P(U_0)$ est défini comme ensemble de $S$-morphismes $U\to X$, il est
clair que $U_0\mapsto P(U_0)$ est un \emph{faisceau d'ensembles}
$\mathscr{P}$ sur $Y_0$. Utilisons ce fait pour prouver que l'application
\[
  h:\Gamma(U_0,\mathscr{G})\times P(U_0)\longrightarrow P(U_0)
\]
définissant la structure de torseur sur $P(U_0)$ est aussi indépendante du
choix de $V$ et $W$, et en outre que, si $U'\subset U$ est un second ouvert
affine de $Y$, $U'_0$ son image réciproque dans $Y_0$, le diagramme
\[
  \xymatrix{
    \Gamma(U_0,\mathscr{G})\times P(U_0)\ar[d]\ar[r]^h
      & P(U_0)\ar[d]\\
    \Gamma(U'_0,\mathscr{G})\times P(U'_0)\ar[r]_{h'}
      & P(U'_0)
  }
  \tag{16.5.14.3}
  \label{IV.16.5.14.3}
\]
est commutatif (les flèches verticales étant les opérateurs de restriction).
En vertu de la remarque précédente, on est ramené à prouver la commutativité
du diagramme précédent lorsque $h$ est défini comme ci-dessus à partir des
ouverts affines $V$, $W$ et $h'$ à partir d'ouverts
